\section{Maximal Rectangle}
\label{sec:sq-maximal-rectangle}

\problemstatement{Given an $n \times m$ binary matrix, find the area of
the largest rectangle that contains only $1$s.}

\begin{patternbox}
\emph{``Largest all-$1$s rectangle in a binary matrix''} is Largest
Rectangle in Histogram (Section~\ref{sec:sq-largest-rectangle-histogram})
applied \textbf{once per row}, where each row defines a histogram of
``how tall a column of consecutive $1$s ends here.'' The keyword to
notice: a 2D ``maximal rectangle of some property'' problem often reduces
to a 1D technique applied row by row, with state carried forward between
rows.
\end{patternbox}

\subsection*{Understanding the problem carefully}

For each row $r$, define $\textit{height}[j]$ as the number of
consecutive $1$s ending at row $r$ in column $j$, looking upward (so
$\textit{height}[j] = 0$ if $\textit{matrix}[r][j] = 0$, and
$\textit{height}[j] = \textit{height}_{\text{previous row}}[j] + 1$ if
$\textit{matrix}[r][j] = 1$). This $\textit{height}$ array is exactly a
histogram: a rectangle of all $1$s that has its \emph{bottom edge} on row
$r$ corresponds exactly to a rectangle in this histogram, since the
histogram's height at column $j$ tells us exactly how far upward a solid
column of $1$s extends from row $r$.

\begin{ideabox}[Run Largest-Rectangle-in-Histogram Once per Row]
Maintain a running $\textit{height}$ array across rows, updating it as
described above after moving to each new row. After updating
$\textit{height}$ for row $r$, run
Algorithm~\ref{sec:sq-largest-rectangle-histogram}'s histogram algorithm
on it, and record the resulting maximum area as a candidate for the
overall answer. After processing every row, the largest area seen across
all rows is the answer.
\end{ideabox}

\subsection*{Why checking every row's histogram captures every possible rectangle}

Every all-$1$s rectangle in the matrix has \emph{some} specific bottom
row $r$. When we process row $r$, the running $\textit{height}$ array at
that point correctly reflects, for every column, how far a solid run of
$1$s extends upward from row $r$ -- so the histogram-area algorithm run on
that exact $\textit{height}$ array is guaranteed to find the largest
rectangle with bottom edge on row $r$ specifically. Since every possible
all-$1$s rectangle has its bottom edge on \emph{some} row, running this
check for every row and taking the overall maximum is guaranteed not to
miss any candidate rectangle.

\subsection*{Algorithm}

\begin{enumerate}
  \item Initialize $\textit{height}[0..m-1] \gets 0$ and
        $\textit{maxArea} \gets 0$.
  \item For each row $r$ from $0$ to $n-1$:
  \begin{enumerate}
    \item For each column $j$: if $\textit{matrix}[r][j] = 1$:
          $\textit{height}[j] \gets \textit{height}[j]+1$; else:
          $\textit{height}[j] \gets 0$.
    \item Run the largest-rectangle-in-histogram algorithm on
          $\textit{height}$; update $\textit{maxArea}$ with the result.
  \end{enumerate}
  \item Return \textit{maxArea}.
\end{enumerate}

\subsection*{Dry run}

\dryrun{Take the matrix
$\begin{bmatrix} 1&0&1&0&0 \\ 1&0&1&1&1 \\ 1&1&1&1&1 \\ 1&0&0&1&0
\end{bmatrix}$.}

\begin{itemize}
  \item Row $0$: $\textit{height}=[1,0,1,0,0]$. Largest histogram
        rectangle: each isolated bar of height $1$, area $1$.
        $\textit{maxArea}=1$.
  \item Row $1$: $\textit{height}=[2,0,2,1,1]$. Largest rectangle: bars
        at indices $2,3,4$ with heights $2,1,1$ give area
        $1\times3=3$ (limited by the shortest, height $1$); also
        indices $0$ alone gives $2$; index $2$ alone gives $2$.
        Best here is $3$. $\textit{maxArea}=3$.
  \item Row $2$: $\textit{height}=[3,1,3,2,2]$. Largest rectangle:
        indices $2,3,4$ with heights $3,2,2$ give area $2\times3=6$.
        $\textit{maxArea}=6$.
  \item Row $3$: $\textit{height}=[4,0,0,3,0]$. Largest rectangle:
        index $0$ alone gives $4$; index $3$ alone gives $3$. Best here
        is $4$, not exceeding $6$.
\end{itemize}

The maximal rectangle has area $\boxed{6}$.

\subsection*{C++ implementation}

\begin{lstlisting}
#include <bits/stdc++.h>
using namespace std;

int largestRectangleArea(vector<int> heights) {
    heights.push_back(0);
    int n = heights.size();
    stack<int> st;
    long long maxArea = 0;

    for (int i = 0; i < n; i++) {
        while (!st.empty() && heights[i] < heights[st.top()]) {
            int j = st.top();
            st.pop();
            int width = st.empty() ? i : i - st.top() - 1;
            maxArea = max(maxArea, (long long)heights[j] * width);
        }
        st.push(i);
    }
    return maxArea;
}

int maximalRectangle(vector<vector<int>>& matrix) {
    int n = matrix.size(), m = matrix[0].size();
    vector<int> height(m, 0);
    int maxArea = 0;

    for (int r = 0; r < n; r++) {
        for (int j = 0; j < m; j++) {
            height[j] = (matrix[r][j] == 1) ? height[j] + 1 : 0;
        }
        maxArea = max(maxArea, largestRectangleArea(height));
    }
    return maxArea;
}

int main() {
    vector<vector<int>> matrix = {
        {1,0,1,0,0},
        {1,0,1,1,1},
        {1,1,1,1,1},
        {1,0,0,1,0}
    };
    cout << "Maximal rectangle area: " << maximalRectangle(matrix) << endl;
    return 0;
}
\end{lstlisting}

\begin{complexitybox}
\textbf{Time Complexity.} Updating the height array for each row costs
$O(m)$, and running the histogram algorithm costs $O(m)$, for a total of
$O(m)$ per row across $n$ rows, giving
\[
  O(nm).
\]
\textbf{Space Complexity.} $O(m)$ for the running height array and the
histogram algorithm's stack.
\end{complexitybox}

\begin{takeawaybox}
A 2D ``maximal region with some property'' problem often decomposes into
``run a 1D algorithm on a derived array, once per row (or column), with
the derived array updated incrementally as you move between rows.'' Once
you have Largest Rectangle in Histogram, Maximal Rectangle requires no new
stack logic at all -- only the row-by-row height bookkeeping is new.
\end{takeawaybox}
